\documentclass[11pt,a4paper]{article}
\usepackage[utf8]{inputenc}
\usepackage[T1]{fontenc}
\usepackage[french]{babel}
\usepackage[margin=2.2cm]{geometry}
\usepackage{booktabs,longtable,array,xcolor,hyperref,enumitem}
\usepackage[htt]{hyphenat}
\hypersetup{colorlinks=true,linkcolor=blue!50!black,urlcolor=blue!50!black}
\newcommand{\code}[1]{\texttt{#1}}
\newcommand{\prio}[1]{\textbf{\textcolor{#1}{\textbullet}}}
\setlist{nosep,leftmargin=*}

\title{VetNutriMP --- Inventaire de la dette technique}
\author{Audit statique du dépôt (branche \code{ccr-757b66d9-sd7476}, commit \code{ef0bc36})}
\date{31 octobre 2026}

\begin{document}
\maketitle

\section*{Méthode}
Analyse statique par comptage (\code{grep}, \code{wc}, \code{md5sum}, \code{git ls-files}) sur
\code{composeApp/src}. Aucune compilation ni exécution des tests n'a été faite. Les
chiffres sont des indicateurs : un \code{catch} ou un \code{!!} n'est pas un défaut en soi,
mais leur densité signale les zones fragiles.

Légende des priorités : \prio{red!80!black}~haute, \prio{orange}~moyenne, \prio{gray}~basse.

\section{Vue d'ensemble chiffrée}
\begin{center}
\begin{tabular}{lr}
\toprule
Indicateur & Valeur \\
\midrule
Fichiers Kotlin (\code{src}, toutes cibles) & 547 \\
Lignes de code Kotlin & \(\approx 136\,600\) \\
dont \code{View/} (101 fichiers) & 61\,329 (\(\approx 45\,\%\)) \\
Lignes de tests (\code{commonTest} + \code{desktopTest}) & \(\approx 10\,500\) \\
Ratio tests / production & \(\approx 7{,}7\,\%\) \\
Fichiers de plus de 1\,000 lignes & 25 \\
\code{catch (e: Exception|Throwable)} dans \code{commonMain} & 480 \\
Opérateurs \code{!!} dans \code{commonMain} & 142 \\
\code{mutableStateOf} dans \code{View/} & 669 \\
Marqueurs \code{TODO}/\code{FIXME} & 16 \\
Taille de \code{.git} & 59\,Mo \\
\bottomrule
\end{tabular}
\end{center}

\section{Architecture}

\subsection*{\prio{red!80!black} A1 --- Composables géants}
La couche \code{View/} concentre près de la moitié du code. Les plus gros fichiers :
\code{SettingsView.kt} (3\,386 l.), \code{AnimalDetailView.kt} (3\,287), \code{ImportUtils.kt} (3\,043),
\code{AnalyseGraphiqueAlimentsView.kt} (3\,041), \code{RationsView.kt} (2\,981),
\code{NewReferenceEvEditView.kt} (2\,695), \code{ConsultationFullScreenEditView.kt} (2\,305),
\code{StartupScreen.kt} (2\,248), \code{MultiNutrientAdjustmentDialog.kt} (2\,231).
\code{AnimalDetailView.kt} ne déclare que 3 fonctions \code{@Composable} pour 3\,287 lignes :
ce sont des blocs monolithiques, difficiles à relire, à prévisualiser et à recomposer
efficacement. Les 669 \code{mutableStateOf} dans les vues indiquent qu'une part importante
de l'état d'écran vit hors des ViewModels.

\emph{Action :} découper par section d'écran (un fichier par onglet ou panneau), en
commençant par les fichiers modifiés le plus souvent.

\subsection*{\prio{red!80!black} A2 --- Violation de la règle « Views $\not\to$ Repository »}
\code{CLAUDE.md} interdit aux Composables de toucher les Repository ; 33 fichiers de
\code{View/} les référencent pourtant, dont \code{RationsView}, \code{FoodSearchComponent},
\code{MultiNutrientAdjustmentDialog}, \code{DetailNutrimentAnalysis}, \code{cardNutrient},
les graphiques de \code{AnalyseGraphique/} et \code{SettingsSections/Recipe*}.
La logique d'accès aux données et une partie des calculs se trouvent donc dans l'UI,
hors de portée des tests unitaires.

\emph{Action :} faire passer ces appels par des méthodes de ViewModel ; traiter en
priorité les vues de calcul (analyse nutritionnelle, ajustement multi-nutriments).

\subsection*{\prio{orange} A3 --- ViewModels « maison » et cycle de vie manuel}
Aucun des 26 ViewModels n'hérite de \code{androidx.lifecycle.ViewModel} alors que la
dépendance est déclarée. Chacun crée son \code{CoroutineScope(AppDispatchers.Main + job)}
et expose un \code{clear()} appelé à la main dans des \code{DisposableEffect} de \code{App.kt}.
Tout oubli d'appel fuit des coroutines. \code{AnimalDetailViewModel} (2\,665 l.),
\code{EquationViewModel} (1\,178) et le couple \code{AnimalViewModel}/\code{AnimalListViewModel}
cumulent trop de responsabilités.

\subsection*{\prio{orange} A4 --- Scopes de coroutines non maîtrisés hors ViewModel}
Scopes globaux sans annulation dans \code{BackupService}, \code{StartupService},
\code{EquationRepository}, \code{LocalizationManager} (avec \code{Dispatchers.Main} direct, en
contournement de \code{AppDispatchers}), \code{BiblioRefRepository:138}
(\code{CoroutineScope(AppDispatchers.IO).launch} \emph{fire-and-forget}) et
\code{DatabaseFoodRepository:67} (scope \emph{Main} dans un Repository).

\subsection*{\prio{orange} A5 --- Règle \code{withContext(AppDispatchers.IO)} non respectée}
Plusieurs Repository ont des \code{suspend fun} sans aucun \code{withContext} :
\code{AnimalRepository}, \code{ConseilRepository}, \code{EquationRepository} (19 fonctions),
\code{ExportImportRepository} (11, 1\,174 l., appelle des DAO directement),
\code{FoodRepository}, \code{HtmlSectionRepository}, \code{PreferencesRepository},
\code{ExamGradingRepository}. Certains déléguent peut-être à des couches qui changent de
dispatcher ; à vérifier au cas par cas.

\subsection*{\prio{gray} A6 --- Code bloquant}
\code{runBlocking} reste présent dans \code{AppDatabase.kt}, \code{desktopMain/Utils/FileUtils.kt}
et \code{windowsMain/main.kt}, plus l'abstraction \code{BlockingRunner} (\code{expect/actual}).
Risque de gel de l'UI s'il est appelé depuis le thread principal.

\section{Robustesse}

\subsection*{\prio{red!80!black} R1 --- Gestion d'erreurs par avalement}
480 blocs \code{catch (e: Exception)} ou \code{Throwable} dans \code{commonMain}, dont au moins
25 vides et une trentaine limités à un \code{println}. Pour une application de calcul
nutritionnel, une erreur silencieuse produit un résultat faux sans signal. Attraper
\code{Exception} dans une coroutine intercepte aussi \code{CancellationException} et casse
l'annulation structurée.

\emph{Action :} journaliser via un logger unique, relancer \code{CancellationException},
remonter les erreurs métier dans l'état UI.

\subsection*{\prio{orange} R2 --- \code{!!} et journalisation}
142 \code{!!} dans \code{commonMain} (crashs potentiels sur données importées incomplètes) ;
14 \code{println} restants alors que \code{slf4j} est présent côté desktop.

\subsection*{\prio{orange} R3 --- Fonctionnalités inachevées (TODO)}
Plusieurs TODO touchent le métier, pas seulement l'UI :
\begin{itemize}
  \item \code{RationCalculator.kt:50,62} : évaluation des équations personnalisées non implémentée ;
  \item \code{NutrientQuantity.kt:24} : conversions entre unités non implémentées ;
  \item \code{MultiNutrientAdjustmentDialog.kt:1426} : ajustement énergétique final absent ;
  \item \code{ExcelFoodService.kt:183} : espèce forcée à 0 ;
  \item \code{SettingsView.kt:3238} : réglage non sauvegardé dans \code{PreferencesStorage} ;
  \item \code{AppDatabase.kt:222,229} : sauvegarde automatique / export JSON non implémentés ;
  \item \code{ConsultationEditDialog.kt:124} : callback \code{onEdit} vide (bouton inopérant).
\end{itemize}

\section{Code mort et exemples en production}
\subsection*{\prio{orange} M1}
\begin{itemize}
  \item \code{Example/} et \code{Examples/} (5 fichiers) compilés dans \code{commonMain} ;
  \item \code{Repository/ConseilRepositoryExample.kt} ;
  \item \code{Greeting.kt} (gabarit de projet KMP) ;
  \item \code{LegacyMigration.kt} (13 l.) à côté de \code{desktopMain/Service/LegacyV2Migration.kt} (1\,604 l.) ;
  \item à la racine : \code{DataB\_Usage\_Example.kt}, \code{TestDataBFiltering.kt}, \code{temp\_histogramme.kt},
        \code{test\_*.kt} (5 fichiers), et un dossier \code{commonMain/} contenant une copie de
        sources \code{koalaplot} et \code{qrgenerator} hors du module Gradle.
\end{itemize}

\section{Build et dépendances}
\subsection*{\prio{red!80!black} B1 --- Incohérences de versions}
\begin{itemize}
  \item \code{koalaplot} : catalogue en 0.9.0, mais \code{build.gradle.kts} code en dur 0.8.0 ;
  \item \code{room} = \code{2.7.0-rc01} (version candidate en production ; \code{CLAUDE.md} indique 2.7.0) ;
  \item \code{kotlinx-coroutines} 1.7.3 avec Kotlin 2.2.20 : très en retard ;
  \item deux versions de \code{kotlinx-datetime} (0.4.1 et 0.6.1) dans le catalogue ;
  \item \code{compose-plugin} 1.6.11, \code{compose} 1.7.8, \code{compose-multiplatform} 1.7.0,
        \code{ui-util} 1.7.0 et \code{material-icons-extended} 1.7.3 codés en dur : versions Compose disparates ;
  \item \code{libs.okio} déclaré deux fois ; \code{junit} 4, \code{junit-jupiter} 5 et \code{testng} au catalogue
        sans usage cohérent ; dépendances commentées laissées en place.
\end{itemize}
\emph{Action :} tout faire passer par \code{libs.versions.toml}, purger les entrées inutilisées.

\subsection*{\prio{orange} B2 --- Données dupliquées}
\code{vetnutri\_export\_init.json} (\(\approx 20\) Mo) existe en 4 exemplaires :
\code{commonMain/resources}, \code{androidMain/assets}, \code{iosApp/iosApp/Resources}
(identiques) et \code{iosApp/iosApp/} (version différente), plus une copie horodatée
\code{\dots 20-06-01-825.json}. Une mise à jour des données doit être reportée à la main
partout ; la copie divergente d'iOS est déjà le symptôme. À générer par une tâche Gradle
depuis une source unique.

\subsection*{\prio{orange} B3 --- Fichiers lourds ou parasites versionnés}
\begin{itemize}
  \item 3 \code{.aab} de release (\(\approx 57\) Mo) toujours suivis malgré la règle
        \code{composeApp/release/} du \code{.gitignore} (ajoutés avant elle) ;
  \item \code{ani.json}, \code{animaux.json}, \code{food.json} (15--19 Mo chacun), \code{oza2.json},
        \code{equations.json}, \code{test\_vbnr.json} à la racine ;
  \item journaux \code{build\_log.txt}, \code{ios\_build.log}, \code{ios\_build\_final.log} ;
        icônes « \code{copie} » en double ; \code{composeApp/bin/} ; \code{iosApp/*/File.txt} ;
  \item \code{.git} pèse 59 Mo : un nettoyage d'historique (\code{git filter-repo}) serait nécessaire
        pour récupérer l'espace, à coordonner car il réécrit l'historique.
\end{itemize}

\section{Tests}
\subsection*{\prio{orange} T1}
\(\approx 50\) fichiers de tests (\(\approx 10\,500\) lignes, 7,7\,\% du code). Le moteur
pur (\code{Data/} : troupeau, exploration multiration, agrégation, analyse) est bien couvert,
ainsi que les migrations Room. En revanche : aucun test sur la couche \code{View/} (45\,\% du code),
un seul test de ViewModel (\code{PlanEvolutifInteractionTest}), aucun sur
\code{ExportImportRepository}, \code{ImportUtils} n'a qu'un test JSON pour 3\,043 lignes. La logique
logée dans les Composables (A2) échappe mécaniquement aux tests.

\section{Documentation}
\subsection*{\prio{gray} D1}
\begin{itemize}
  \item \code{CLAUDE.md} est désynchronisé : 51 vues annoncées (101 réelles), 23 ViewModels (26),
        8 fichiers \code{DataBase/} (9), 24 migrations (25 objets \code{Migration}) ;
  \item 15 fichiers \code{.md} à la racine (\code{ANALYSE\_*}, \code{OPTIMISATION\_*}, \code{README\_*},
        \code{RESUME\_*}) dont certains se recoupent ; à déplacer dans \code{docs/} ou à archiver ;
  \item fichier \code{Task} et \code{docs/Untitled.json} sans rôle identifiable.
\end{itemize}

\section{Plan de remboursement proposé}
\begin{longtable}{p{0.5cm}p{9.5cm}p{2cm}p{2.3cm}}
\toprule
 & Action & Effort & Gain \\
\midrule
1 & Purger fichiers parasites et \code{.aab} de l'index, supprimer code d'exemple et fichiers \code{test\_*} racine & faible & lisibilité, taille \\
2 & Unifier les versions dans \code{libs.versions.toml}, sortir Room de la RC, monter coroutines & faible--moyen & stabilité build \\
3 & Source unique pour \code{vetnutri\_export\_init.json} générée au build & moyen & cohérence données \\
4 & Politique d'erreurs : logger commun, \code{CancellationException} relancée, suppression des \code{catch} vides & moyen & fiabilité calculs \\
5 & Terminer ou retirer les TODO métier (équations personnalisées, conversions d'unités) & moyen & exactitude \\
6 & Ajouter \code{withContext(IO)} aux Repository concernés, rattacher les scopes orphelins & moyen & réactivité UI \\
7 & Sortir les accès Repository des 33 vues vers les ViewModels, puis tester ces ViewModels & élevé & testabilité \\
8 & Découper les 9 fichiers de plus de 2\,000 lignes & élevé & maintenabilité \\
\bottomrule
\end{longtable}

\section{État après corrections (branche \code{ccr-757b66d9-sd7476})}
Chaque étape a été validée par la CI GitHub (tests JVM common + desktop, tests unitaires
Android, APK debug). La compilation iOS n'est lancée que manuellement et n'a pas été exécutée.

\begin{longtable}{p{0.5cm}p{3.2cm}p{11cm}}
\toprule
 & Statut & Détail \\
\midrule
1 & fait & \code{.aab}, journaux, fichiers \code{test\_*}/temporaires, JSON de travail, copie de
           koalaplot, icônes « copie » retirés ; code d'exemple et 5 ViewModels jamais instanciés
           supprimés. L'historique Git n'a \emph{pas} été réécrit (les 59\,Mo restent dans
           \code{.git}). \\
2 & fait & Catalogue unique ; Room 2.7.0 ; coroutines 1.10.2 ; entrées mortes retirées. \\
3 & fait & Source unique ; assets Android générés au build ; copie iOS resynchronisée (elle
           avait divergé) et vérifiée par \code{InitialDataSyncTest}. \\
4 & fait & \code{AppLog}, \code{rethrowIfCancellation()} dans 688 \code{catch}, 88 \code{catch}
           vides journalisés, 156 \code{println}/\code{printStackTrace} remplacés. \\
5 & fait & TODO résolus ou code mort supprimé (\code{RationCalculator},
           \code{NutrientQuantityExt}, \code{createDatabaseBackup} qui renvoyait \code{true}
           sans rien faire) ; sélecteurs de référentiels et création depuis un modèle de conseil
           implémentés. Plus aucun \code{TODO} dans le code. \\
6 & fait & \code{withContext(IO)} dans 6 Repository ; scopes nommés avec \code{SupervisorJob} ;
           références bibliographiques fictives insérées dans les bases vides supprimées. \\
7 & partiel & Conseils, écran de démarrage et export d'animal passent par des ViewModels.
           \textbf{Préférences} : 12 instances de \code{PreferencesRepository} construites dans les vues
           (fiche animal, rations, détail de nutriment, carte d'analyse, Paramètres) remplacées par
           l'instance d'\code{AppContainer}, exposée en \code{StateFlow} par
           \code{AnimalDetailViewModel}, \code{MultiRationViewModel} et \code{SettingsViewModel} :
           une préférence enregistrée dans Paramètres est vue sans rechargement (test
           \code{PreferencesRepositorySharedStateTest}) ; 3 blocs de chargement morts supprimés.
           Reste : \code{equationRepository} transmis aux fonctions de calcul du domaine
           (\code{getEnergie}, \code{getNutrientWithComplementary}) --- à remplacer par une
           interface de résolution d'équations dédiée. \\
8 & fait & Plus aucun fichier de \code{View/} au-delà de 2\,000 lignes ; découpage vérifié par
           captures (400\,×\,800, 800\,×\,1000, 1400\,×\,900).
           \code{AnimalDetailView} 3\,287\,$\to$\,1\,343 l. (\code{WideScreenLayout} et
           \code{NarrowScreenLayout}, identiques à 92\,\%, fusionnés en \code{AnimalDetailLayout} ;
           export dans \code{AnimalDetailExport.kt}) ;
           \code{SettingsView} 3\,386\,$\to$\,1\,381 (dialogues de sélection d'export, onglet Préférences) ;
           \code{AnalyseGraphiqueAlimentsView} 3\,041\,$\to$\,1\,041 (graphiques, données) ;
           \code{NewReferenceEvEditView} 2\,695\,$\to$\,608 (nutriments, coefficients, onglets
           Équations/Compléments) ; \code{RationsView} 2\,981\,$\to$\,1\,947 ;
           \code{ConsultationFullScreenEditView} 2\,305\,$\to$\,1\,403 ;
           \code{StartupScreen} 2\,248\,$\to$\,1\,716 ;
           \code{MultiNutrientAdjustmentDialog} 2\,231\,$\to$\,925 : l'algorithme d'ajustement
           (1\,265 l., sans UI) passe dans \code{Data/AjustementMultiNutriments.kt}, ce qui supprime
           les imports \code{Data} $\to$ \code{View.AnalNut} (\code{RationConstraintAdjuster},
           \code{ExplorationMultiration}) et \code{ViewModel} $\to$ \code{View}.
           Les composables principaux restent longs (\code{RationsView} $\approx$ 1\,880 l.,
           \code{StartupScreen} $\approx$ 1\,700, éditeur de consultation $\approx$ 1\,400) : les
           découper demande d'extraire leur état. \\
-- & partiel & \code{!!} : 104 restants dans \code{commonMain} (142 à l'audit). Revue des sites
           exécutés après un clic ou une attente : un seul non gardé, corrigé
           (\code{besoinEnergetiqueStandard!!} dans l'ajustement multi-nutriments, plantage tant que le BEE
           n'est pas calculé) ; un autre réécrit sans \code{!!}. Les autres suivent une garde
           (\code{contains}, \code{isEmpty}, test de nullité sur un état Compose délégué, qui empêche le
           smart-cast). Cycle de vie manuel des ViewModels (\code{clear()} depuis \code{App.kt})
           inchangé : \code{MultiRationViewModel} et \code{SettingsViewModel} héritent déjà de
           \code{ViewModel}, les autres non. \\
\bottomrule
\end{longtable}

\section{Interface et petits écrans (vérification visuelle)}
Rendu réel de l'application desktop, hors écran, avec le pilote
\code{GuideScreenshotDriver} (nouvelle variable \code{VETNUTRI\_GUIDE\_SIZE}) en
400\,×\,800 (téléphone), 800\,×\,1000 (tablette) et 1400\,×\,900 (bureau), sur une base
jetable : accueil, ajout d'animal, aliments, données de calcul, paramètres, fiche animal,
consultations, rations.

\begin{longtable}{p{3.2cm}p{11.5cm}}
\toprule
Écran & Défaut constaté $\to$ correction \\
\midrule
Consultations & Liste de 160\,px et détail côte à côte sur téléphone (« Consultation s ») $\to$
  une seule colonne sous 600\,dp : liste, puis détail en plein écran
  (\code{AppSizes.breakpointCompact}). \\
Détail de consultation & Coefficients K1--K5 sur deux colonnes écrasés lettre par lettre
  (« N on dé fi ni ») $\to$ une colonne quand la place manque ; bloc dupliqué factorisé
  (\code{CoefficientsConsultation}). \\
Paramètres, Calculs et 7 autres barres d'onglets & Libellés coupés en milieu de mot
  (« Équatio ns », « Inter face ») $\to$ \code{AdaptiveTabRow} : onglets défilants quand ils
  ne tiennent pas. \\
Liste des aliments & Filtres tronqués (« Sélection ner », « Tout es »), indication de recherche
  coupée en hauteur $\to$ un filtre par ligne sur téléphone, indication sur une ligne avec
  « … ». \\
Rations & Badge « APPORTS » coupé en « APP/ORT/S » $\to$ en-tête de la carte d'analyse en
  \code{FlowRow} ; grille collée aux bords $\to$ marges ; grille de cartes vides sans
  consultation $\to$ seul le message de sélection reste. \\
Ajout d'animal & Listes Espèce/Sexe à 280\,dp avec double bordure, champ Race à bordure
  pâle (air désactivé), boutons de date empilés $\to$ \code{ComboBox} pleine largeur,
  bordure standard, boutons côte à côte. \\
Fiche animal & « Date de naissance :2026-10-31 » collés ; icône d'Ordonnance identique à
  Paramètres $\to$ colonne d'étiquettes élargie et atténuée ; icône document. \\
Éditeur de consultation & Flèche retour annoncée « Rechercher » (accessibilité) $\to$
  « Retour ». \\
\bottomrule
\end{longtable}

Suites traitées :
\begin{itemize}
  \item \textbf{Couleur secondaire} : \code{\#80A9A9} avec texte noir (boutons Retour, Partager,
        Import rapide, \textminus/+ à l'air désactivé) remplacé par \code{\#4E7D7D} avec texte
        blanc (contraste $\approx$ 4,7:1). Trois boutons qui forçaient déjà du blanc deviennent lisibles.
  \item \textbf{Traductions} : \code{androidMain/assets/strings\_*.json} (1 clé manquante) supprimées,
        générées au build ; copies iOS \code{strings\_fr/en/zh.json} (10 clés manquantes : clés brutes
        affichées) resynchronisées par \code{syncInitialDataIos} et vérifiées par
        \code{InitialDataSyncTest}.
  \item Libellés français sans accents (sélection d'export : « Selection des equations »,
        « toggle tout », « Selectionnes ») corrigés ; barre d'actions de l'Ordonnance sur téléphone :
        ordonnance pleine largeur, libellés du compte rendu sur deux lignes au lieu d'être tronqués.
  \item Quatre fichiers Kotlin vides et suivis supprimés (\code{BackupViewModel},
        \code{TestApiEnvelope}, \code{FullscreenToggle}, \code{GenericGraphComponent}).
  \item Filtres de la sélection d'aliments (Graphique Aliments) : un par ligne sur téléphone.
  \item Fichiers morts retirés : \code{AnimalDetailView.kt.old}, \code{RationsView.kt.bak2},
        \code{View/oza2.json}, \code{View/vetfood.json} (1,9\,Mo), \code{PreferencesView.kt} (jamais appelé).
\end{itemize}


\end{document}
